%% ===== Abstract =====
% [27/10] Viết lại cả abstract sau review: số theo phương án B (Compilation 71%, O-e và O-h 4%);
% bỏ vế "only in 2026" (V5); "relies on human judgment" thành "9 of the 12 ... include a human check"
% (comment 1: chỉ 3/9 bài LIB có người chấm mà không chạy code); thêm pilot, 52% đến 90% (P1);
% tách câu dài để mỗi dấu câu chỉ một chức năng.

% [29/10] "Nine" thành "Ten" theo Bảng 10 (Google migrations mã lại); "six gaps" thành "five gaps" theo Discussion (G1--G5).
\begin{abstract}
\cmt{Claude}{[R01, major, LaTeX] Review of the 29/10 version: blue [Claude: R..] notes are review comments; the full report with the reasoning and the checks behind each one is Review\_Claude\_Seminar\_2\_2910.md at the repository root. Before the clean build: switching to reviewfalse in macros.tex currently fails to compile (``Undefined control sequence'' for hcrev, first in Section 5.1), because hcrev and revtag are defined only in the review branch; add empty definitions to the else branch (details in macros.tex).}
\new{Large language models (LLMs) are now used to translate code between languages and to migrate it across versions, libraries and frameworks. Whether the generated code behaves like the original is decided by an oracle, and the choice of oracle can change the reported result. We present a systematic mapping study of how behavioral equivalence is evaluated in LLM-based code translation and migration, covering the main proceedings of ICSE, FSE and ASE from 2024 to 2026. \cmt{Claude}{[R02] ``Main proceedings'' here also covers Companion-volume papers (some from workshops), posters, doctoral symposium and journal-first presentations (Section 4.1); most readers will take it to mean the research-track volume. See R10. The abstract is also long (about 310 words) for the EASE version.} From 4,451 papers we selected 49 and coded each from its full text. The most common oracles are compilation (71\%) and existing tests (57\%); formal verification, LLM-generated tests and LLM judges are rare (4\% each). Nearly a fifth of the papers use no oracle that runs the code. Oracle use differs by context: none of the 5 papers with a legacy source uses existing tests, and 9 of the 12 on library migration include a human check. Thirty-three papers report two or more oracles on the same outputs. In the 18 that give both a compile rate and a behavioral pass rate, behavior trails compilation by 3 to 98 percentage points. \cmt{Claude}{[R03] The 98 comes from the DepTrans compile-only-reward ablation (see R21). Adding the median gap (19 pp) would keep the range from being read as typical.} Ten of the 49 papers state no limitation of their own oracle. A pilot then varies only the oracle. We judged an LLM's translations of 164 HumanEval-X tasks from Python to Java with four oracles: compilation, the benchmark tests, CodeBLEU and a second LLM as judge. Depending on the oracle and its threshold, 52\% to 90\% of the same translations count as correct. The judge accepts half of those that fail the tests and even rejects 41\% of the benchmark's own reference solutions. \cmt{Claude}{[R04, major] The lower end, 52\%, uses CodeBLEU at 0.7, a threshold the pilot does not treat as main; at the main threshold the range is 67\% to 90\%. Consider ``67\% to 90\% (52\% under a stricter similarity threshold)''. Also say the judge is small (Llama 3.1 8B, 4-bit), as the Conclusion does, so that 50\% and 41\% are not read as properties of LLM judges in general.} A success rate is thus meaningful only together with the oracle that produced it. From these results we derive five gaps between evaluation practice and real migration, along with a checklist for reporting oracle provenance. We release the data and the pilot's harness.}
\end{abstract}
